%% ma: L11-B16-0094
%% lop: 11
%% chuong: 5
%% bai: 16
%% dang: 11.16.1
%% loai: TL
%% muc: VD
%% dap_an: "-"
%% nguon: "Ảnh giáo viên gửi 25/11/2026: Giới hạn hàm số dạng 0/0 (3 bài, có ĐS), Bài 3 câu 2"
%% kien_thuc: "Lớp 11 Bài 16 (giới hạn của hàm số); hằng đẳng thức x^k-1=(x-1)(x^(k-1)+...+1); tổng 1+2+...+n (Lớp 11 Bài 6 cấp số cộng)"
%% trang_thai: chinh-thuc
\begin{ex}
Tính giới hạn $A=\lim\limits_{x\to1}\dfrac{x^{50}-1}{x^2-3x+2}$.
\loigiai{$\dfrac{(x-1)(x^{49}+\dots+x+1)}{(x-1)(x-2)}=\dfrac{x^{49}+\dots+x+1}{x-2}$. Vậy $A=\dfrac{50}{-1}=-50$.}
\end{ex}
